\documentclass[11pt,a4paper]{article}

% ---------- Packages ----------
\usepackage[utf8]{inputenc}
\usepackage[T1]{fontenc}
\usepackage[margin=2.5cm]{geometry}
\usepackage{amsmath,amssymb,amsthm}
\usepackage{mathtools}
\usepackage{hyperref}
\usepackage{xcolor}
\usepackage{listings}
\usepackage{algorithm}
\usepackage{algpseudocode}
\usepackage{enumitem}
\usepackage{booktabs}
\usepackage{tikz}
\usepackage{fancyvrb}

\hypersetup{
  colorlinks=true,
  linkcolor=blue!60!black,
  urlcolor=blue!60!black,
  citecolor=blue!60!black,
}

% ---------- Theorem environments ----------
\newtheorem{theorem}{Theorem}
\newtheorem{lemma}[theorem]{Lemma}
\newtheorem{proposition}[theorem]{Proposition}
\newtheorem{corollary}[theorem]{Corollary}
\theoremstyle{definition}
\newtheorem{definition}[theorem]{Definition}
\newtheorem{remark}[theorem]{Remark}
\newtheorem{example}[theorem]{Example}

% ---------- Go listing style ----------
\definecolor{gocomment}{RGB}{106,153,85}
\definecolor{gokeyword}{RGB}{0,102,204}
\definecolor{gostring}{RGB}{163,21,21}
\definecolor{gonumber}{RGB}{181,90,40}
\definecolor{gobg}{RGB}{248,248,248}

\lstdefinelanguage{Go}{
  keywords={break,default,func,interface,select,case,defer,go,map,struct,
    chan,else,goto,package,switch,const,fallthrough,if,range,type,
    continue,for,import,return,var,true,false,nil,iota},
  keywordstyle=\color{gokeyword}\bfseries,
  ndkeywords={string,bool,int,int8,int16,int32,int64,uint,uint8,uint16,
    uint32,uint64,uintptr,byte,rune,float32,float64,complex64,complex128,
    error,big,Int,Element},
  ndkeywordstyle=\color{gokeyword!70!black},
  sensitive=true,
  comment=[l]{//},
  morecomment=[s]{/*}{*/},
  commentstyle=\color{gocomment}\itshape,
  stringstyle=\color{gostring},
  morestring=[b]",
  morestring=[b]`,
}

\lstset{
  language=Go,
  basicstyle=\ttfamily\small,
  backgroundcolor=\color{gobg},
  numbers=left,
  numberstyle=\tiny\color{gray},
  numbersep=6pt,
  frame=single,
  rulecolor=\color{gray!50},
  breaklines=true,
  showstringspaces=false,
  tabsize=4,
  captionpos=b,
  xleftmargin=1em,
  framexleftmargin=1em,
}

% ---------- Shortcuts ----------
\newcommand{\F}{\mathbb{F}}
\newcommand{\Z}{\mathbb{Z}}
\newcommand{\N}{\mathbb{N}}
\DeclareMathOperator{\sk}{sk}
\DeclareMathOperator{\vk}{vk}
\DeclareMathOperator{\QR}{QR}

\title{\bfseries Constraint-Friendly Map-to-Curve Relations\\
\large From the GMMZ Construction to the y-Increment Fix}
\author{A Technical Companion}
\date{\today}

\begin{document}
\maketitle

\begin{abstract}
This document presents a unified technical treatment of two recent papers on
map-to-elliptic-curve-group relations for constraint programming. First, we
describe the GMMZ construction (Groth, Malvai, Miller, Zhang, Asiacrypt~2025),
which replaces the expensive inner cryptographic hash by a nondeterministic
map-to-curve relation whose witness is verified by a handful of algebraic
constraints. Second, we describe the algebraic attack of El~Housni and Bünz,
which exploits the order-3 automorphism of $j=0$ curves (such as BN254) to
produce a concrete BLS signature forgery whenever $M \cdot T \gtrsim \sqrt{q}$.
Third, we present their countermeasure, the \emph{y-increment} construction,
which neutralizes the automorphism by encoding messages into the $y$-coordinate
instead of the $x$-coordinate, and simultaneously lifts the field restriction
$q \equiv 3 \pmod 4$. For each construction we provide the mathematical
specification, pseudocode, and a self-contained Go implementation.
\end{abstract}

\tableofcontents

% ============================================================
\section{Background and Notation}
% ============================================================

\subsection{Elliptic curves}

Let $q > 3$ be prime and let $E$ be an elliptic curve over $\F_q$ in short
Weierstrass form
\[
  E : y^2 = f(x), \qquad f(x) = x^3 + ax + b.
\]
The set $E(\F_q) = \{(x,y) \in \F_q^2 : y^2 = f(x)\} \cup \{\mathcal{O}\}$
forms a finite abelian group, where $\mathcal{O}$ is the point at infinity.
The inverse of a point is $-(x,y) = (x,-y)$. A cryptographic subgroup
$G \subseteq E(\F_q)$ of prime order $p$ is singled out; the cofactor
$h = |E(\F_q)|/p$ is typically small (e.g.\ $h = 1$ for BN254, secp256k1).

The \emph{$j$-invariant} classifies a curve up to $\overline{\F_q}$-isomorphism:
\[
  j(E) = 1728 \cdot \frac{4a^3}{4a^3 + 27 b^2}.
\]
Curves with $j=0$ (i.e.\ $a=0$, $y^2 = x^3+b$) admit an extra order-3
automorphism
\[
  \phi : (x,y) \mapsto (\omega x, y),
\]
where $\omega \in \F_q$ is a primitive cube root of unity (existing whenever
$q \equiv 1 \pmod 3$). Crucially, $\phi(P) = [\lambda] P$ for a scalar
$\lambda$ with $\lambda^2 + \lambda + 1 \equiv 0 \pmod p$. This structural
property is what powers the attack of Section~\ref{sec:attack}.

\subsection{Constraint programming for zero-knowledge proofs}

In constraint programming one specifies a relation $R \subseteq X \times W$
between public input $x \in X$ and witness $w \in W$ through algebraic
equations. A prover demonstrates that it knows $w$ with $(x,w) \in R$ by
producing a succinct proof. The cost of proving grows linearly with the
number of constraints, so every saved multiplication matters.

Hash-to-curve functions are used for multiset hashing (for zkVM memory
consistency) and for BLS signatures (for proof-of-stake aggregation). The
classical construction uses an inner cryptographic hash (SHA-256, Poseidon,
MiMC) followed by a deterministic map-to-curve. The inner hash dominates
the constraint cost: SHA-256 alone costs $\sim 7000$ constraints per
invocation in R1CS, whereas a complete map-to-curve relation can be verified
in only a few dozen constraints.

% ============================================================
\section{The GMMZ x-Increment Construction}
\label{sec:gmmz}
% ============================================================

\subsection{Core idea}

GMMZ observes that in the Generic Group Model (GGM), the encoding of a group
element already randomizes discrete logarithms. There is therefore no need
for an inner cryptographic hash, provided the message space $M$ is small
enough that the collision probability $q_{\text{oracle}} \cdot |M| / p$
remains negligible. The map-to-group step alone, expressed as a
nondeterministic relation, suffices.

The construction interprets a message $m \in M \subseteq \F_q$ directly as
(the base of) an $x$-coordinate. Since a random $x$ lies on the curve only
with probability $\approx 1/(2h)$, an iterative \emph{increment-and-check}
is used: the prover tries $x = m \cdot T + k$ for $k = 0,1,\dots,T-1$ and
supplies the successful $k$ as witness.

\subsection{The relation}

Fix an elliptic curve $E : y^2 = x^3 + ax + b$ over $\F_q$ with $q \equiv 3
\pmod 4$, a subgroup $G \subseteq E(\F_q)$ of prime order $p$, a message
space $M \subseteq \F_q$, and a tweak bound $T \in \N$ with $M \cdot T < q$.

\begin{definition}[GMMZ x-increment relation]
\label{def:x-increment}
The relation $R_x \subseteq M \times G \times W$ consists of all triples
$(m,(x,y),w)$ with $w = (k,z)$ satisfying
\begin{align}
  & m \in M, \quad k \in [0,T), \label{eq:x1}\\
  & x = m \cdot T + k, \label{eq:x2}\\
  & z^2 = y, \label{eq:x3}\\
  & y^2 = x^3 + a x + b, \label{eq:x4}\\
  & (x,y) \in G. \label{eq:x5}
\end{align}
\end{definition}

\paragraph{Role of each constraint.}
\begin{itemize}
  \item \textbf{\eqref{eq:x2} Injectivity.} The multiplier $T$ forces disjoint
  $x$-ranges $[mT, (m{+}1)T)$ for distinct $m$; given $x$ one recovers
  $(m,k)$ uniquely by Euclidean division. The side condition $M \cdot T < q$
  prevents modular wrap-around.

  \item \textbf{\eqref{eq:x3} Inverse exclusion.} On an elliptic curve,
  $(x,y)$ and $(x,-y)$ are both valid points but inverse to each other.
  Without a canonical choice, two witnesses would map one message to two
  points. The constraint $z^2 = y$ forces $y$ to be a quadratic residue. When
  $q \equiv 3 \pmod 4$, $-1$ is a quadratic non-residue, so exactly one of
  $\{y, -y\}$ is a QR—canonicality is obtained for free.

  \item \textbf{\eqref{eq:x4},\eqref{eq:x5} Curve and subgroup membership.}
  These ensure $(x,y)$ is actually an element of $G$.
\end{itemize}

\subsection{The constructor (prover side)}

\begin{algorithm}[H]
\caption{GMMZ x-Increment Constructor}
\label{alg:x-inc-construct}
\begin{algorithmic}[1]
\Require curve $E : y^2 = x^3+ax+b$ over $\F_q$ with $q \equiv 3 \pmod 4$,
subgroup $G$, message $m \in M$, tweak bound $T$.
\State $k \gets 0$
\While{$k < T$}
  \State $x \gets m \cdot T + k$
  \State $f \gets x^3 + a x + b$
  \If{$f^{(q-1)/2} = 1$ in $\F_q$}
     \Comment{quadratic residue test}
     \State $y \gets f^{(q+1)/4}$
     \State $z \gets y^{(q+1)/4}$
        \Comment{$z^2 = y$ since $y$ is a QR}
     \If{$(x,y) \in G$}
        \State \Return $((x,y),\,(k,z))$
     \EndIf
  \EndIf
  \State $k \gets k+1$
\EndWhile
\State \Return $\bot$ \Comment{failure probability $\approx 2^{-T}$}
\end{algorithmic}
\end{algorithm}

\subsection{Applications}

\subsubsection*{Multiset hashing.}
Given an injective and inverse-excluding map-to-group relation, define
\[
  \mathrm{MSH}(S) = \sum_{m \in S} \mathrm{MapToGroup}(m).
\]
Incrementality is immediate: adding $m$ to the multiset adds $(x_m,y_m)$ to
the digest. Collision resistance holds in EC-GGM with probability at most
$(2 \sum_m |S_m^\pm| \cdot Q + 2Q^2)/p$ for $Q$ oracle queries.

\subsubsection*{BLS signatures.}
\begin{itemize}
  \item $\mathrm{KeyGen}$: sample $\sk \leftarrow \Z_p$; set
  $\vk = [\sk] g_2 \in G_2$.
  \item $\mathrm{Sign}(m)$: compute $(h_m, w_m) \gets
  \mathrm{Construct}(m)$; output $\sigma = [\sk] h_m$.
  \item $\mathrm{Verify}(\vk, m, \sigma, w_m)$: check $(m, h_m, w_m) \in
  R_{\text{M2G}}$ and $e(\sigma, g_2) = e(h_m, \vk)$.
\end{itemize}

\subsection{Go implementation}
\label{subsec:go-xinc}

We implement the constructor for BN254. The curve is
$y^2 = x^3 + 3$ over $\F_q$ with $q \equiv 3 \pmod 4$, so the square root
and QR test are a single exponentiation each.

\begin{lstlisting}[caption={GMMZ x-increment constructor (Go, BN254).},
label={lst:x-inc-go}]
package mapcurve

import (
    "errors"
    "math/big"

    bn254 "github.com/consensys/gnark-crypto/ecc/bn254"
    "github.com/consensys/gnark-crypto/ecc/bn254/fp"
)

// TweakBound is the maximum number of tweak attempts.
const TweakBound = 256

// XIncWitness is the witness (k, z) for the x-increment relation.
type XIncWitness struct {
    K uint32   // tweak in [0, T)
    Z fp.Element // square root of y: z^2 = y
}

// XIncConstruct maps a message m (given as a field element) to a BN254 G1
// point by the GMMZ x-increment procedure.
//
//   x = m*T + k,    y = sqrt(f(x)),   z = sqrt(y)
//
// Returns the point (x,y) and the witness (k,z). The failure probability
// after T = 256 trials is about 2^{-256}.
func XIncConstruct(m *fp.Element) (bn254.G1Affine, XIncWitness, error) {
    var (
        tBig = big.NewInt(int64(TweakBound))
        mBig = new(big.Int)
        base fp.Element // base = m * T  (as a field element)
    )
    m.BigInt(mBig)
    base.SetBigInt(new(big.Int).Mul(mBig, tBig))

    var (
        x, y, f  fp.Element
        kElement fp.Element
    )

    for k := uint32(0); k < TweakBound; k++ {
        // x = m*T + k
        kElement.SetUint64(uint64(k))
        x.Add(&base, &kElement)

        // f = x^3 + ax + b  (for BN254: a=0, b=3)
        computeRHS(&f, &x)

        // QR test: f is a QR iff f^{(q-1)/2} = 1
        if !isQR(&f) {
            continue
        }

        // y = sqrt(f) = f^{(q+1)/4}  (valid because q = 3 mod 4)
        y = sqrtModP3Mod4(&f)

        // z = sqrt(y) = y^{(q+1)/4}  (valid because y is a QR by choice)
        z := sqrtModP3Mod4(&y)

        // Build the affine point and verify subgroup membership.
        var P bn254.G1Affine
        P.X = x
        P.Y = y
        if !P.IsOnCurve() || !P.IsInSubGroup() {
            continue
        }

        return P, XIncWitness{K: k, Z: z}, nil
    }
    return bn254.G1Affine{}, XIncWitness{},
        errors.New("x-increment: no valid tweak found")
}

// computeRHS sets out = x^3 + 3  for BN254 (a=0, b=3).
func computeRHS(out, x *fp.Element) {
    var x2, x3 fp.Element
    x2.Square(x)
    x3.Mul(&x2, x)
    var three fp.Element
    three.SetUint64(3)
    out.Add(&x3, &three)
}

// isQR returns true iff f is a quadratic residue in F_q.
func isQR(f *fp.Element) bool {
    // Legendre symbol: f^{(q-1)/2} = +-1 for nonzero f.
    return f.Legendre() == 1
}

// sqrtModP3Mod4 returns f^{(q+1)/4} when q = 3 mod 4.
func sqrtModP3Mod4(f *fp.Element) fp.Element {
    var out fp.Element
    // gnark-crypto exposes a direct Sqrt; we use it here for brevity.
    out.Sqrt(f)
    return out
}
\end{lstlisting}

\begin{lstlisting}[caption={Verifier-side algebraic check of the
x-increment relation.}, label={lst:x-inc-verify}]
// XIncVerify re-checks the five algebraic conditions of Definition 1.
// In a SNARK this is compiled to a handful of R1CS constraints.
func XIncVerify(m *fp.Element, P *bn254.G1Affine, w *XIncWitness) bool {
    if w.K >= TweakBound {
        return false
    }

    // Condition (x1): x = m*T + k
    var (
        tElem, kElem, base, x fp.Element
    )
    tElem.SetUint64(uint64(TweakBound))
    kElem.SetUint64(uint64(w.K))
    base.Mul(m, &tElem)
    x.Add(&base, &kElem)
    if !x.Equal(&P.X) {
        return false
    }

    // Condition (x3): z^2 = y  (canonical QR root of y)
    var z2 fp.Element
    z2.Square(&w.Z)
    if !z2.Equal(&P.Y) {
        return false
    }

    // Conditions (x4)-(x5): (x,y) on curve and in subgroup.
    return P.IsOnCurve() && P.IsInSubGroup()
}
\end{lstlisting}

% ============================================================
\section{The Algebraic Attack}
\label{sec:attack}
% ============================================================

\subsection{The automorphism blind spot}

The EC-GGM of Groth–Shoup models the encoding $\pi : \Z_p \to G$ as a random
bijection that is identity-preserving and inverse-preserving, but it does
not account for the order-3 automorphism $\phi$ that exists on every $j=0$
curve over a field with $q \equiv 1 \pmod 3$. The adversary can compute
$\phi(P)$ from coordinates alone (one field multiplication: $x \mapsto \omega
x$), without issuing any oracle query—yet the security proof only bounds
the advantage in terms of oracle queries. This is the model gap.

\subsection{Lattice argument}

Given $\phi(P_1) = P_2$, the $x$-coordinates satisfy
\[
  x_2 \equiv \omega \cdot x_1 \pmod q.
\]
The set of all such integer pairs is the lattice
\[
  L = \{(a,b) \in \Z^2 \,:\, b \equiv \omega\cdot a \pmod q\},
  \qquad \det L = q.
\]
Minkowski's theorem guarantees a nonzero $(a,b) \in L$ with
\[
  \max(|a|,|b|) \le \sqrt{q}.
\]
When $M \cdot T > \sqrt{q}$, both $|a|$ and $|b|$ fit inside the encoding
range $[0, MT)$. The short vector is computed in $O(\log^2 q)$ time by the
extended Euclidean algorithm on $(q, \omega)$, or in $O(1)$ closed form for
parameterized families.

\begin{proposition}[BN short vector, El Housni--Bünz]
For BN curves parameterized by seed $u$ with $q(u) = 36u^4 + 36u^3 + 24u^2 +
6u + 1$, the pair
\[
  a(u) = 6u^2 + 2u + 1, \qquad b(u) = 2u
\]
satisfies $\omega(u) \cdot a(u) \equiv b(u) \pmod{q(u)}$ with
$\omega(u) = 18u^3 + 18u^2 + 9u + 1$, and $\max(|a|,|b|) \le \sqrt{q}$.
\end{proposition}

\subsection{The BLS forgery}

\begin{algorithm}[H]
\caption{Existential BLS signature forgery on a $j=0$ curve}
\label{alg:forgery}
\begin{algorithmic}[1]
\Require public key $\vk$, signing oracle $\mathcal{O}_{\sk}$, BN seed $u$.
\State Compute $(a,b)$ via the Cornacchia factorization (or
Proposition above for BN curves).
\State $x_1 \gets |a| \bmod q$; let $y_1$ be the QR root of $x_1^3 + b$.
\State Let $m_1 \gets \lfloor x_1 / T \rfloor$.
\State Recompute $P_1 = H(m_1)$ by running the public x-increment
constructor (the oracle sees the same point by determinism).
\State Query $\sigma_1 \gets \mathcal{O}_{\sk}(m_1)$.
\State \Comment{Apply $\phi^2$ as a pure coordinate map on both points.}
\State Set $P_2 \gets (\omega^2 \cdot x(P_1),\; y(P_1))$.
\State Set $\sigma_2 \gets (\omega^2 \cdot x(\sigma_1),\; y(\sigma_1))$.
\State Let $m_2 \gets \lfloor x(P_2) / T \rfloor$.
\State \Return $(m_2, \sigma_2, P_2)$.
\end{algorithmic}
\end{algorithm}

The forgery bypasses the scalar $\lambda$ entirely: since $\phi$ is a group
endomorphism, $\phi^2(\sigma_1) = \phi^2([\sk] P_1) = [\sk]\phi^2(P_1) =
[\sk] P_2$ purely by group-action linearity. Verification succeeds because
\[
  e(\sigma_2, g_2)
  = e([\sk]\phi^2(P_1), g_2)
  = e(P_2, [\sk] g_2)
  = e(P_2, \vk).
\]
Importantly, computing $\phi^2$ is a single field multiplication ($x \mapsto
\omega^2 x$) — the adversary never invokes the group addition oracle, which
is why the EC-GGM security argument misses this attack entirely.

\subsection{Go implementation of the attack}

\begin{lstlisting}[caption={BN254 algebraic attack: building the short
vector and forging a signature.}, label={lst:attack-go}]
package attack

import (
    "fmt"
    "math/big"

    bn254 "github.com/consensys/gnark-crypto/ecc/bn254"
    "github.com/consensys/gnark-crypto/ecc/bn254/fp"
)

// BN254Seed is the curve-defining seed u (also called x_0). The BN254
// moduli are q(u) = 36u^4 + 36u^3 + 24u^2 + 6u + 1, p(u) = q(u) - 6u^2.
var BN254Seed, _ = new(big.Int).SetString(
    "4965661367192848881", 10)

// ComputeShortVector returns (a(u), b(u), omega(u)) from Proposition 1.
//   a(u) = 6u^2 + 2u + 1
//   b(u) = 2u
//   omega(u) = 18u^3 + 18u^2 + 9u + 1   (a primitive cube root of unity in Fq)
func ComputeShortVector(u *big.Int) (a, b, omega *big.Int) {
    u2 := new(big.Int).Mul(u, u)
    u3 := new(big.Int).Mul(u2, u)

    // a = 6u^2 + 2u + 1
    a = new(big.Int).Mul(big.NewInt(6), u2)
    a.Add(a, new(big.Int).Mul(big.NewInt(2), u))
    a.Add(a, big.NewInt(1))

    // b = 2u
    b = new(big.Int).Mul(big.NewInt(2), u)

    // omega = 18u^3 + 18u^2 + 9u + 1
    omega = new(big.Int).Mul(big.NewInt(18), u3)
    omega.Add(omega, new(big.Int).Mul(big.NewInt(18), u2))
    omega.Add(omega, new(big.Int).Mul(big.NewInt(9), u))
    omega.Add(omega, big.NewInt(1))

    return
}

// ForgeBLS produces a valid BLS signature on a fresh message m2 using a
// single signing-oracle query on a crafted m1. Exploits the order-3
// automorphism of BN254 (a j=0 curve).
//
// The trick: phi(x,y) = (omega*x, y) is an automorphism with phi(P) =
// [lambda]P on G1, so applying phi^2 to sigma1 = [sk]*H(m1) as a pure
// coordinate map yields phi^2(sigma1) = [sk]*phi^2(H(m1)) by group-action
// linearity. We never need to know lambda — phi is a coordinate map.
//
// Returns (m2, sigma2, P2), where P2 = phi^2(H(m1)) is the forged hash
// point the verifier will reconstruct. The caller then runs the standard
// BLS pairing check e(sigma2, g2) =? e(P2, vk).
func ForgeBLS(
    signOracle func(m fp.Element) bn254.G1Affine,
) (m2 fp.Element, sigma2 bn254.G1Affine, P2 bn254.G1Affine, err error) {

    aInt, _, omegaInt := ComputeShortVector(BN254Seed)

    // Take x1 = |a| mod q, split into (m1, k1_unused) by x1 = m1*T + k.
    qModulus := fp.Modulus()
    x1Int := new(big.Int).Mod(aInt, qModulus)
    T := big.NewInt(int64(TweakBound))
    m1Int := new(big.Int)
    k1Unused := new(big.Int)
    m1Int.DivMod(x1Int, T, k1Unused)

    var m1 fp.Element
    m1.SetBigInt(m1Int)

    // Recompute H(m1) ourselves (public deterministic function); the oracle
    // produces the same point, so we can read its coordinates.
    H1, _, cerr := XIncConstruct(&m1)
    if cerr != nil {
        return fp.Element{}, bn254.G1Affine{}, bn254.G1Affine{},
            fmt.Errorf("XIncConstruct(m1) failed: %w", cerr)
    }

    // Query the signing oracle.
    sigma1 := signOracle(m1)

    // Apply phi^2 as coordinate map: phi^2(x, y) = (omega^2 * x, y).
    var omegaFp, omega2 fp.Element
    omegaFp.SetBigInt(omegaInt)
    omega2.Square(&omegaFp)

    // Forged hash point: P2 = phi^2(H(m1)).
    P2.X.Mul(&omega2, &H1.X)
    P2.Y = H1.Y
    if !P2.IsOnCurve() {
        return fp.Element{}, bn254.G1Affine{}, bn254.G1Affine{},
            fmt.Errorf("P2 off curve (impossible on j=0)")
    }

    // Forged signature: sigma2 = phi^2(sigma1) = [sk] * phi^2(H(m1)).
    sigma2.X.Mul(&omega2, &sigma1.X)
    sigma2.Y = sigma1.Y

    // Extract m2 from P2.x by Euclidean division.
    var x2Big big.Int
    P2.X.BigInt(&x2Big)
    var m2Big, k2Big big.Int
    m2Big.DivMod(&x2Big, T, &k2Big)
    m2.SetBigInt(&m2Big)

    return m2, sigma2, P2, nil
}
\end{lstlisting}

% ============================================================
\section{The y-Increment Countermeasure}
\label{sec:yinc}
% ============================================================

\subsection{Core idea}

Since $\phi(x,y) = (\omega x, y)$ \emph{fixes} the $y$-coordinate, encoding
the message into $y$ rather than $x$ makes the automorphism
\emph{same-message}: $\phi(P)$ encodes the same $m$ as $P$, so no cross-message
collision is possible. Two bonus properties fall out:
\begin{itemize}
  \item \textbf{Free inverse exclusion.} The inverse point $-P = (x,-y)$ has
  $y$-coordinate $q - y \notin [0, MT)$ whenever $M \cdot T < q/2$. The QR
  witness chain $z^2 = y$ is no longer needed.
  \item \textbf{No field restriction.} Because inverse exclusion is
  geometric, the condition $q \equiv 3 \pmod 4$ disappears entirely.
\end{itemize}

\subsection{The relation}

\begin{definition}[y-increment relation]
\label{def:y-increment}
Let $E : y^2 = x^3 + a x + b$ be an elliptic curve over $\F_q$ with prime
subgroup $G$, message space $M \subseteq \F_q$, and tweak bound $T$
satisfying $M \cdot T < \sqrt{q}$. The relation $R_y \subseteq M \times G
\times W$ consists of all $(m,(x,y),w)$ with $w = k$ satisfying
\begin{align}
  & m \in M, \quad k \in [0,T),\\
  & y = m \cdot T + k,\\
  & y^2 = x^3 + a x + b,\\
  & (x,y) \in G.
\end{align}
\end{definition}

\paragraph{Why $M \cdot T < \sqrt{q}$?}
The tightened bound reflects the true security threshold of
\emph{any} such relation in EC-GGM: the $Q^2/p$ birthday term in GMMZ's own
theorem caps security at $\sqrt p$, which for a prime-order subgroup of a
curve over $\F_q$ is essentially $\sqrt q$.

\subsection{The constructor}

Given $y$, recovering $x$ requires solving the cubic $x^3 + a x + c = 0$
with $c = b - y^2$. Two cases:

\paragraph{$j = 0$ curves ($a = 0$).}
The cubic degenerates to $x^3 = y^2 - b$, i.e.\ a cube root. When
$q \equiv 1 \pmod 3$, exactly $1/3$ of field elements are cubes, and a
candidate $y$ succeeds with probability $\approx 1/3$. With $T = 256$, the
failure probability is $(2/3)^{256} < 2^{-150}$.

\paragraph{General curves ($j \ne 0, 1728$).}
Solve the depressed cubic $x^3 + ax + c = 0$ by Cardano's formula, with a
Legendre-symbol branch on the discriminant
\[
  \Delta = -4a^3 - 27 c^2.
\]

\begin{algorithm}[H]
\caption{y-Increment Constructor (uniform over $j=0$ and $j \ne 0$)}
\label{alg:y-inc-construct}
\begin{algorithmic}[1]
\Require curve $E : y^2 = x^3 + a x + b$ over $\F_q$, subgroup $G$,
message $m \in M$, tweak bound $T$.
\State $k \gets 0$
\While{$k < T$}
  \State $y \gets m \cdot T + k$
  \State $c \gets b - y^2$
  \If{$a = 0$}
     \Comment{$j=0$ case: cube root}
     \State $x \gets (-c)^{1/3} \bmod q$ if a cube root exists
  \Else
     \Comment{general case: Cardano}
     \State $x \gets $ root of $X^3 + aX + c$ via Cardano
  \EndIf
  \If{$x \ne \bot$ and $(x,y) \in G$}
     \State \Return $((x,y), k)$
  \EndIf
  \State $k \gets k+1$
\EndWhile
\State \Return $\bot$
\end{algorithmic}
\end{algorithm}

\subsection{Security statement}

\begin{theorem}[Informal, El Housni--Bünz]
\label{thm:y-sec}
Let $E$ be a $j=0$ curve over $\F_q$ with $M \cdot T < \sqrt q$. In the
standard Groth–Shoup EC-GGM, any adversary making at most $Q$ oracle queries
against the y-increment BLS scheme has forging advantage at most
\[
  \mathrm{Adv} \le \frac{c \cdot M T \cdot Q + Q^2}{p}
\]
for an absolute constant $c$. In particular, security reaches the DLP
threshold $\min(p/(MT), \sqrt p)$ \emph{without} any model augmentation.
\end{theorem}

\noindent Conceptually, the proof proceeds exactly like GMMZ's Theorem~3,
but now $\phi$ is harmless because it acts as the identity on the encoding
coordinate $y$. No extra oracle and no extra lattice argument are needed.

\subsection{Go implementation}

\subsubsection*{$j=0$ constructor (BN254).}

\begin{lstlisting}[caption={y-increment constructor for BN254 ($j=0$).},
label={lst:y-inc-go}]
package mapcurve

import (
    "errors"
    "math/big"
    "sync"

    bn254 "github.com/consensys/gnark-crypto/ecc/bn254"
    "github.com/consensys/gnark-crypto/ecc/bn254/fp"
)

// YIncWitness is the witness k for the y-increment relation.
type YIncWitness struct {
    K uint32
}

// YIncConstructBN254 maps m to a BN254 G1 point using the y-increment
// procedure. For j=0 curves, x-recovery reduces to a cube root.
func YIncConstructBN254(m *fp.Element) (bn254.G1Affine, YIncWitness, error) {
    var (
        tBig = big.NewInt(int64(TweakBound))
        mBig = new(big.Int)
        base fp.Element
    )
    m.BigInt(mBig)
    base.SetBigInt(new(big.Int).Mul(mBig, tBig))

    var (
        y, y2, negC, x, kElem fp.Element
    )
    var bFp fp.Element
    bFp.SetUint64(3) // BN254: b = 3

    for k := uint32(0); k < TweakBound; k++ {
        kElem.SetUint64(uint64(k))
        y.Add(&base, &kElem)

        // negC = y^2 - b ; we want x^3 = negC
        y2.Square(&y)
        negC.Sub(&y2, &bFp)

        // Cube root in Fq (BN254 has q = 1 mod 3, so only 1/3 of elements
        // are cubes). If negC is not a cube, skip this tweak.
        if !cubeRootInFq(&x, &negC) {
            continue
        }

        var P bn254.G1Affine
        P.X = x
        P.Y = y
        if !P.IsOnCurve() || !P.IsInSubGroup() {
            continue
        }
        return P, YIncWitness{K: k}, nil
    }
    return bn254.G1Affine{}, YIncWitness{},
        errors.New("y-increment: no valid tweak found")
}

// cubeRootInFq computes out = t^{1/3} in BN254's F_q, or returns false if
// t is not a cube.
//
// BN254's q - 1 = 2 * 3^2 * m with gcd(m, 6) = 1 (i.e. v_3(q-1) = 2).
// Therefore 9 | q - 1, and t^{(q+8)/27} cubed equals t * zeta^k for some k,
// where zeta is a primitive 9th root of unity — NOT a mere cube root of
// unity. So the correction loop iterates up to 9 times, not 3.
//
// All arithmetic stays inside gnark-crypto's Montgomery form via Exp/Mul.
var (
    cubeExpOnce  sync.Once
    cubeExpE     *big.Int   // (q+8)/27
    cubeResidueE *big.Int   // (q-1)/3
    zeta9        fp.Element // primitive 9th root of unity in F_q
)

func cubeExpInit() {
    cubeExpOnce.Do(func() {
        q := fp.Modulus()
        cubeExpE = new(big.Int).Add(q, big.NewInt(8))
        cubeExpE.Div(cubeExpE, big.NewInt(27))
        cubeResidueE = new(big.Int).Sub(q, big.NewInt(1))
        cubeResidueE.Div(cubeResidueE, big.NewInt(3))

        // Build zeta9: pick some small g that is NOT a 9th-power residue
        // and set zeta9 = g^{(q-1)/9}.  Require order exactly 9 (so its
        // cube != 1).
        qm1Over9 := new(big.Int).Sub(q, big.NewInt(1))
        qm1Over9.Div(qm1Over9, big.NewInt(9))

        var one, g, pow9, cubed fp.Element
        one.SetOne()
        for i := int64(2); i < 1000; i++ {
            g.SetUint64(uint64(i))
            pow9.Exp(g, qm1Over9)
            if pow9.Equal(&one) {
                continue
            }
            cubed.Square(&pow9)
            cubed.Mul(&cubed, &pow9)
            if !cubed.Equal(&one) {
                zeta9 = pow9
                return
            }
        }
        panic("cubeRoot: failed to find a primitive 9th root of unity")
    })
}

func cubeRootInFq(out, t *fp.Element) bool {
    if t.IsZero() {
        out.SetZero()
        return true
    }
    cubeExpInit()

    // Cubic-residue test: t^{(q-1)/3} == 1.
    var r, one fp.Element
    r.Exp(*t, cubeResidueE)
    one.SetOne()
    if !r.Equal(&one) {
        return false
    }

    // Candidate: cand = t^{(q+8)/27}. Then cand^3 = t * zeta9^k for some k.
    var cand, cube fp.Element
    cand.Exp(*t, cubeExpE)

    for i := 0; i < 9; i++ {
        cube.Square(&cand)
        cube.Mul(&cube, &cand)
        if cube.Equal(t) {
            *out = cand
            return true
        }
        cand.Mul(&cand, &zeta9)
    }
    return false
}
\end{lstlisting}

\subsubsection*{Verifier.}

\begin{lstlisting}[caption={y-increment verifier: only four simple
checks, no residuosity chain.}, label={lst:y-inc-verify}]
// YIncVerify re-checks the four algebraic conditions of Definition 3.
// In an R1CS, this compiles to ~16 constraints.
func YIncVerify(m *fp.Element, P *bn254.G1Affine, w *YIncWitness) bool {
    if w.K >= TweakBound {
        return false
    }

    // y = m*T + k
    var tElem, kElem, base, y fp.Element
    tElem.SetUint64(uint64(TweakBound))
    kElem.SetUint64(uint64(w.K))
    base.Mul(m, &tElem)
    y.Add(&base, &kElem)
    if !y.Equal(&P.Y) {
        return false
    }

    return P.IsOnCurve() && P.IsInSubGroup()
}
\end{lstlisting}

\subsubsection*{General case: Cardano solver for $j \ne 0, 1728$.}

\begin{lstlisting}[caption={Cardano's depressed-cubic solver over $\F_q$.
This is what the y-increment needs when $a \ne 0$.},
label={lst:cardano}]
// SolveDepressedCubic tries to find x in F_q with x^3 + a*x + c = 0.
// Returns (x, true) on success, or (_, false) if no F_q-root exists.
//
// Strategy:
//   D^2 = c^2/4 + a^3/27                 (Cardano discriminant term)
//   If D is a square in F_q:
//       u = (-c/2 + D)^{1/3} in F_q
//       x = u - a/(3u)
//   Else:
//       u = (-c/2 + D)^{1/3} in F_{q^2}, and the real root is its trace.
//
// For the y-increment with j != 0, we typically only need the "D-square"
// branch for simplicity. Full generality requires F_{q^2} arithmetic
// (omitted here).
func SolveDepressedCubic(a, c *fp.Element) (fp.Element, bool) {
    var (
        two, three, four, twentySeven fp.Element
        halfC, a3Over27, Dsq          fp.Element
        half                          fp.Element
    )
    two.SetUint64(2)
    three.SetUint64(3)
    four.SetUint64(4)
    twentySeven.SetUint64(27)

    half.Inverse(&two)

    // halfC = c/2
    halfC.Mul(c, &half)

    // a3Over27 = a^3 / 27
    var a2, a3 fp.Element
    a2.Square(a)
    a3.Mul(&a2, a)
    var inv27 fp.Element
    inv27.Inverse(&twentySeven)
    a3Over27.Mul(&a3, &inv27)

    // Dsq = (c/2)^2 + a^3/27
    var halfC2 fp.Element
    halfC2.Square(&halfC)
    Dsq.Add(&halfC2, &a3Over27)

    // Only handle the case Dsq is a square in F_q.
    if Dsq.Legendre() != 1 {
        return fp.Element{}, false
    }

    var D fp.Element
    D.Sqrt(&Dsq)

    // inside = -c/2 + D
    var inside fp.Element
    inside.Neg(&halfC)
    inside.Add(&inside, &D)

    var u fp.Element
    if ok := cubeRootInFq(&u, &inside); !ok {
        return fp.Element{}, false
    }

    if u.IsZero() {
        // Degenerate case; root is x = -a/(3u) undefined; use Cardano's
        // depressed limit x = (3c/a) if a != 0.
        if a.IsZero() {
            return fp.Element{}, false
        }
        var x fp.Element
        x.Mul(&three, c)
        var aInv fp.Element
        aInv.Inverse(a)
        x.Mul(&x, &aInv)
        return x, true
    }

    // x = u - a/(3u)
    var threeU, invThreeU, aOver3u, x fp.Element
    threeU.Mul(&three, &u)
    invThreeU.Inverse(&threeU)
    aOver3u.Mul(a, &invThreeU)
    x.Sub(&u, &aOver3u)
    return x, true
}
\end{lstlisting}

% ============================================================
\section{End-to-End Demo}
\label{sec:demo}
% ============================================================

The Go file bundles a \texttt{main} function that exercises all three
pieces end-to-end on BN254. It (i) runs the x-increment construct and
verify on a random 100-bit message, (ii) does the same for the y-increment,
and (iii) spins up a toy BLS signing oracle (the key is sampled fresh with
\texttt{crypto/rand}), asks it for a single signature on a carefully
chosen $m_1$, derives a forged signature on a fresh $m_2$, and checks the
pairing equation to confirm the forgery is accepted.

\begin{lstlisting}[caption={Signing-oracle harness that wraps the
x-increment construction into a BLS signer / verifier.},
label={lst:oracle}]
// BLSOracle is a toy signing oracle: on input m, it runs the x-increment
// constructor to get H(m), then returns [sk] * H(m).
type BLSOracle struct {
    sk      big.Int // secret scalar in [1, p)
    vk      bn254.G2Affine
    g2      bn254.G2Affine
    queried []fp.Element // log of oracle queries
}

func NewBLSOracle() (*BLSOracle, error) {
    _, _, _, g2 := bn254.Generators()

    skInt, err := rand.Int(rand.Reader, fr.Modulus())
    if err != nil {
        return nil, err
    }
    if skInt.Sign() == 0 {
        skInt.SetUint64(1)
    }
    var vk bn254.G2Affine
    vk.ScalarMultiplication(&g2, skInt)
    return &BLSOracle{sk: *skInt, vk: vk, g2: g2}, nil
}

// Sign implements the signing-oracle interface expected by ForgeBLS.
func (o *BLSOracle) Sign(m fp.Element) bn254.G1Affine {
    Hm, _, err := XIncConstruct(&m)
    if err != nil {
        panic(fmt.Sprintf("signing oracle: %v", err))
    }
    o.queried = append(o.queried, m)
    var sigma bn254.G1Affine
    sigma.ScalarMultiplication(&Hm, &o.sk)
    return sigma
}

// pairingCheck returns true iff e(a, b) == e(c, d).
func pairingCheck(a bn254.G1Affine, b bn254.G2Affine,
    c bn254.G1Affine, d bn254.G2Affine) bool {
    var negC bn254.G1Affine
    negC.Neg(&c)
    ok, err := bn254.PairingCheck(
        []bn254.G1Affine{a, negC},
        []bn254.G2Affine{b, d},
    )
    if err != nil {
        return false
    }
    return ok
}
\end{lstlisting}

\begin{lstlisting}[caption={Three demo routines and the \texttt{main}
driver. Each sub-demo is self-contained.}, label={lst:main}]
func demoXIncrement() error {
    banner("1) GMMZ x-increment: construct + verify")
    m, err := sampleMessage(100)
    if err != nil {
        return err
    }
    P, w, err := XIncConstruct(&m)
    if err != nil {
        return fmt.Errorf("construct failed: %w", err)
    }
    if !XIncVerify(&m, &P, &w) {
        return errors.New("verifier rejected a valid witness")
    }
    fmt.Println("  verifier accepts")
    return nil
}

func demoYIncrement() error {
    banner("2) y-increment countermeasure: construct + verify (BN254)")
    m, err := sampleMessage(100)
    if err != nil {
        return err
    }
    P, w, err := YIncConstructBN254(&m)
    if err != nil {
        return fmt.Errorf("construct failed: %w", err)
    }
    if !YIncVerify(&m, &P, &w) {
        return errors.New("verifier rejected a valid witness")
    }
    fmt.Println("  verifier accepts")
    return nil
}

func demoAttack() error {
    banner("3) Algebraic attack: BLS signature forgery on x-increment")
    oracle, err := NewBLSOracle()
    if err != nil {
        return err
    }

    // ForgeBLS now returns the forged hash point P2 alongside (m2, sigma2),
    // so no manual reconstruction is needed.
    m2, sigma2, P2, err := ForgeBLS(oracle.Sign)
    if err != nil {
        return fmt.Errorf("forgery failed: %w", err)
    }
    if oracle.AlreadyQueried(m2) {
        return errors.New(
            "ForgeBLS returned the same message it queried — not a real forgery")
    }

    if !oracle.VerifyForgedPoint(P2, sigma2) {
        return errors.New("pairing check on forgery FAILED")
    }
    fmt.Println("  pairing check accepts the forgery — BLS is broken")
    return nil
}

func main() {
    fmt.Println("map-to-curve on BN254: demo of GMMZ x-increment,")
    fmt.Println("the El Housni-Bunz forgery, and the y-increment countermeasure.")

    if err := demoXIncrement(); err != nil {
        fmt.Printf("x-increment demo FAILED: %v\n", err)
    }
    if err := demoYIncrement(); err != nil {
        fmt.Printf("y-increment demo FAILED: %v\n", err)
    }
    if err := demoAttack(); err != nil {
        fmt.Printf("attack demo FAILED: %v\n", err)
    }
    banner("done")
}
\end{lstlisting}

\subsection*{Running the demo}

Expected output (abbreviated), with the concrete forged points elided:

\begin{Verbatim}[frame=single,fontsize=\small]
map-to-curve on BN254: demo of GMMZ x-increment,
the El Housni-Bunz forgery, and the y-increment countermeasure.

============================================================
  1) GMMZ x-increment: construct + verify
============================================================
  message m = 836...104
  H(m).x    = 213...672
  witness   = (k=1, z=...)
  verifier accepts

============================================================
  2) y-increment countermeasure: construct + verify (BN254)
============================================================
  message m = 475...290
  H(m).y    = 121...648
  witness   = (k=0)
  verifier accepts

============================================================
  3) Algebraic attack: BLS signature forgery on x-increment
============================================================
  forged m2 = 123...
  m2 is distinct from every queried m1
  pairing check accepts the forgery — BLS is broken

============================================================
  done
============================================================
\end{Verbatim}

The key observation in line 3 of the output: the verifier accepts a signature
on a message the signing oracle was never queried on. This is an existential
forgery, the definitional failure of EUF-CMA security.

% ============================================================
\section{Comparison and Takeaways}
% ============================================================

\begin{center}
\begin{tabular}{lcc}
\toprule
 & \textbf{x-increment (GMMZ)} & \textbf{y-increment (EH-B)} \\
\midrule
Encoding coordinate      & $x$ & $y$ \\
Inverse exclusion        & QR witness $z^2 = y$ & Free (range argument) \\
Field restriction        & $q \equiv 3 \pmod 4$ & None \\
Vulnerable to $j=0$ automorphism & \textbf{Yes} & No (fixed by $y$) \\
Security threshold on $M \cdot T$ & Was claimed $\ll q$; true is $\sqrt q$ & $\sqrt q$ \\
R1CS constraints (native, BN254)  & 18 & 16 \\
R1CS constraints (emulated, BLS12-377) & 7741 & 1065 \\
\bottomrule
\end{tabular}
\end{center}

\paragraph{Practical conclusion.}
Hashing into a prime-order subgroup of an elliptic curve can indeed be made
dramatically constraint-friendly by dropping the inner cryptographic hash.
The GMMZ approach correctly captures the main idea, but its security analysis
missed the order-3 automorphism on deployed $j=0$ curves (BN254, BLS12-381,
BLS12-377, Grumpkin), resulting in concrete forgeries whenever $M \cdot T
\gtrsim \sqrt q$. The y-increment fix is both simpler (fewer constraints,
no field restriction) and provably secure in the standard EC-GGM, and should
be preferred in any deployment.

\paragraph{Where to go from here.}
For the curious reader: (i) extend the Go implementations above into a full
BLS signing/verification pipeline, (ii) compile the verifier of
Definition~\ref{def:y-increment} as a gnark circuit and measure constraint
counts, (iii) explore the $j = 1728$ case, where the duality
remark at the end of \S\ref{sec:yinc} shows that x-increment is safe but
y-increment is \emph{not}---an interesting reversal of roles.

\end{document}
